\begin{theorem}[Literal character insertions and the actual closed graph contraction]
    \label{thm:peps_regular_charge_bond_contraction}
    \lean{TNLean.PEPS.regularChargeMatrix_eq_leftRegular_diagonal,
      TNLean.PEPS.graphInsertedBondNetwork_averagingSite_charge}
    \leanok
    \uses{def:peps_regular_charge_coordinates,
      thm:peps_graph_inserted_averaging_bond_state}
    Let $U_x$ denote the left-regular matrix and
    $D_{\chi,p}=\operatorname{diag}(g\mapsto\chi(pg))$.
    Independent translations at the two endpoints give exactly
    \[
        U_xD_{\chi,p}U_{y^{-1}}=A_\chi(p;x,y).
    \]
    Consequently, for a finite simple graph with the regular averaging
    tensor at every vertex, inserting $D_{\chi_e,p_e}$ on every edge
    gives the actual regrouped closed coefficient
    \[
        \Psi(\beta)
          =|G|^{-|V|}\sum_{q:V\to G}
            \prod_{e\in E}
              A_{\chi_e}(p_e;q_{\mathrm{head}(e)},q_{\mathrm{tail}(e)})
                 _{\beta_{e,\mathrm{head}},\beta_{e,\mathrm{tail}}}.
    \]
    This is an auxiliary finite-graph coefficient identity for
    \cite[charge definition and detection]{Schuch2010PEPS}, lines 2449--2486.
    It makes no original-spin measurement or general semiregular
    isometry claim for the character insertions.
\end{theorem}
\begin{proof}\leanok
    The regular permutation matrices leave one surviving intermediate label,
    giving the displayed matrix entry. Substitute this identity into the
    proved closed contraction with arbitrary bond insertions.
\end{proof}

\begin{theorem}[Conjugation of a character insertion by a flux]
    \label{thm:peps_regular_charge_flux_conjugation}
    \lean{TNLean.PEPS.leftRegular_diagonal_character_conjugate}
    \leanok
    \uses{thm:peps_regular_charge_bond_contraction}
    For every coefficient function $\chi:G\to\mathbb C$ and $p,k\in G$,
    \[
        U_kD_{\chi,p}U_{k^{-1}}=D_{\chi,pk^{-1}}.
    \]
    This is the virtual coefficient identity in
    \cite[braiding of charges with fluxes]{Schuch2010PEPS}, lines 2569--2581.
    The physical braid remains a separate assertion.
\end{theorem}
\begin{proof}\leanok
    Evaluate the translated charge matrix with both endpoint translations
    equal to $k$. Its entries vanish off the diagonal; the diagonal entry
    at $g$ is $\chi(pk^{-1}g)$.
\end{proof}
